\section{Introduction}
\label{sec:introduction}

Physics-based simulation has become indispensable for training and evaluating
robot perception and manipulation systems. Platforms such as NVIDIA Isaac
Sim~\cite{Mittal2025Isaac,Mittal2023Orbit,Makoviychuk2021Isaac} leverage
Universal Scene Description (USD) as a scene representation and GPU-accelerated
ray tracing to produce photorealistic synthetic data at scale. Recent efforts
demonstrate that millions of high-fidelity rendered images can train competitive
object detectors~\cite{Singh2024Synthetica} and instance segmentation
models~\cite{Ng2023SynTable} without any real-world labels, making the
rendering pipeline a critical component of modern sim-to-real transfer.

A key contributor to the visual realism of Isaac Sim is NVIDIA's Material
Definition Language (MDL), which provides physically based shading with
multi-layer effects, subsurface scattering, and procedural noise functions.
However, MDL materials carry practical costs: they require the NVIDIA Omniverse
runtime, increase scene load times, elevate GPU memory consumption, and produce
assets incompatible with standard 3D tools and web viewers that expect
UsdPreviewSurface or PBR glTF materials. Practitioners who generate synthetic
training data at scale~\cite{Singh2024Synthetica,Ng2023SynTable} or deploy
digital twins in resource-constrained settings therefore face a recurring
question: \emph{can we simplify MDL materials to UsdPreviewSurface without
degrading the downstream utility of the rendered images?}

Despite the widespread use of Isaac Sim in the robotics and computer vision
communities, no prior work has systematically quantified the consequences of
this material simplification. The sim-to-real gap literature focuses on
the boundary between simulated and real images~\cite{Tobin2017Domain,Tremblay2018Training,Zakharov2022Photo}, but largely ignores the
\emph{within-simulation} rendering gap introduced by material representation
choices. This gap is particularly insidious because it operates silently inside
the data generation pipeline: a practitioner may convert assets for efficiency,
unaware of any impact on the training signal.

In this paper we address this gap with four contributions:
\begin{enumerate}
    \item We present a systematic multi-level quantification of
    MDL-to-UsdPreviewSurface material simplification trade-offs in NVIDIA Isaac
    Sim, measuring visual quality degradation across multiple scenes and
    viewpoints.
    \item We conduct a multi-level evaluation spanning pixel-level image quality
    (PSNR, SSIM~\cite{Wang2004Image}, LPIPS~\cite{Zhang2018Unreasonable}),
    semantic feature preservation (CLIP~\cite{Radford2021Learning} and
    DINOv2~\cite{Oquab2023DINOv2} cosine similarity), and downstream AI task
    performance (object detection confidence).
    \item We derive practical decision guidelines for synthetic data
    practitioners: under what conditions material simplification is safe, and
    when full-fidelity MDL rendering remains necessary.
    \item We release \textbf{ConvertAsset}\footnote{Repository URL withheld for double-blind review.}, an open-source toolkit that
    automates the MDL-to-UsdPreviewSurface conversion pipeline with
    composition-preserving USD traversal, texture remapping, and optional mesh
    simplification and GLB export.
\end{enumerate}
